%  LaTeX support: latex@mdpi.com 
%  For support, please attach all files needed for compiling as well as the log file, and specify your operating system, LaTeX version, and LaTeX editor.

%=================================================================
\documentclass[article,submit,pdftex,moreauthors,mathematics]{Definitions/mdpi}
%\documentclass[preprints,article,submit,pdftex,moreauthors]{Definitions/mdpi} 
% For posting an early version of this manuscript as a preprint, you may use "preprints" as the journal. Changing "submit" to "accept" before posting will remove line numbers.
 
% Below journals will use APA reference format:
% admsci, aieduc, behavsci, businesses, econometrics, economies, education, ejihpe, famsci, games, humans, ijcs, ijfs, jintelligence, journalmedia, jrfm, jsam, languages, peacestud, psycholint, publications, tourismhosp, youth

% Below journals will use Chicago reference format:
% arts, genealogy, histories, humanities, laws, literature, religions, risks, socsci

%--------------------
% Class Options:
%--------------------
%----------
% journal
%----------
% Choose between the following MDPI journals:
% accountaudit, acoustics, actuators, addictions, adhesives, admsci, adolescents, aerobiology, aerospace, agriculture, agriengineering, agrochemicals, agronomy, ai, aichem, aieduc, aieng, aimater, aimed, aipa, air, aisens, algorithms, allergies, alloys, amh, analog, analytica, analytics, anatomia, anesthres, animals, antibiotics, antibodies, antioxidants, applbiosci, appliedchem, appliedmath, appliedphys, applmech, applmicrobiol, applnano, applsci, aquacj, architecture, arm, arthropoda, arts, asc, asi, astronautics, astronomy, atmosphere, atoms, audiolres, automation, axioms, bacteria, batteries, bdcc, behavsci, beverages, biochem, bioengineering, biologics, biology, biomass, biomechanics, biomed, biomedicines, biomedinformatics, biomimetics, biomolecules, biophysica, bioresourbioprod, biosensors, biosphere, biotech, birds, blockchains, bloods, blsf, brainsci, breath, buildings, businesses, cancers, carbon, cardio, cardiogenetics, cardiovascmed, catalysts, cells, ceramics, challenges, chemengineering, chemistry, chemosensors, chemproc, children, chips, cimb, civileng, cleantechnol, climate, clinbioenerg, clinpract, clockssleep, cmd, cmtr, coasts, coatings, colloids, colorants, commodities, complexities, complications, compounds, computation, computers, condensedmatter, conservation, constrmater, cosmetics, covid, crops, cryo, cryptography, crystals, csmf, ctn, culture, curroncol, cyber, dairy, data, ddc, dentistry, dermato, dermatopathology, designs, devices, dhi, diabetology, diagnostics, dietetics, digital, disabilities, diseases, diversity, dna, drones, dynamics, earth, ebj, ecm, ecologies, econometrics, economies, edm, education, eesp, ejihpe, electricity, electrochem, electronicmat, electronics, encyclopedia, endocrines, energies, eng, engproc, entomology, entropy, environments, environremediat, epidemiologia, epigenomes, esa, est, famsci, fermentation, fibers, fintech, fire, fishes, fluids, foods, forecasting, forensicsci, forests, fossstud, foundations, fractalfract, fuels, future, futureinternet, futurepharmacol, futurephys, futuretransp, galaxies, games, gases, gastroent, gastrointestdisord, gastronomy, gels, genealogy, genes, geographies, geohazards, geomatics, geometry, geosciences, geotechnics, geriatrics, germs, glacies, grasses, green, greenhealth, gucdd, hardware, hazardousmatters, healthcare, hearts, hemato, hematolrep, hep, heritage, higheredu, highthroughput, histories, horticulturae, hospitals, humanities, humans, hydrobiology, hydrogen, hydrology, hydropower, hygiene, idr, iic, ijcs, ijem, ijerph, ijfs, ijgi, ijmd, ijms, ijns, ijom, ijpb, ijt, ijtm, ijtpp, ime, immuno, informatics, information, infrastructures, inorganics, insects, instruments, inventions, iot, j, jaestheticmed, jal, jcdd, jcm, jcp, jcrm, jcs, jcto, jdad, jdb, jdream, jemr, jeta, jfb, jfmk, jgbg, jgg, jimaging, jintelligence, jlpea, jmahp, jmmp, jmms, jmp, jmse, jne, jnt, jof, joi, joitmc, joma, jop, joptm, jor, journalmedia, jox, jpbi, jphytomed, jpm, jrfm, jsam, jsan, jtaer, jvd, jzbg, kidneydial, kinasesphosphatases, knowledge, labmed, laboratories, lae, land, languages, laws, life, lights, limnolrev, lipidology, liquids, literature, livers, logics, logistics, lubricants, lymphatics, machines, macromol, magnetism, magnetochemistry, make, marinedrugs, materials, materproc, mathematics, mca, measurements, medicina, medicines, medsci, membranes, merits, metabolites, metals, meteorology, methane, metrics, metrology, micro, microarrays, microbiolres, microelectronics, micromachines, microorganisms, microplastics, microwave, minerals, mining, mmphys, modelling, molbank, molecules, mps, msf, mti, multimedia, muscles, nanoenergyadv, nanomanufacturing, nanomaterials, ncrna, ndt, network, neuroglia, neuroimaging, neurolint, neurosci, nitrogen, notspecified, nursrep, nutraceuticals, nutrients, obesities, occuphealth, oceans, ohbm, onco, optics, oral, organics, organoids, osteology, oxygen, pandemics, parasites, parasitologia, particles, pathogens, pathophysiology, peacestud, pediatrrep, pets, pharmaceuticals, pharmaceutics, pharmacoepidemiology, pharmacy, philosophies, photochem, photonics, phycology, physchem, physics, physiologia, plants, plasma, platforms, pollutants, polymers, polysaccharides, populations, poultry, powders, precipitation, precisoncol, preprints, proceedings, processes, prosthesis, proteomes, psf, psychiatryint, psychoactives, psycholint, publications, purification, quantumrep, quaternary, qubs, radiation, rdt, reactions, realestate, receptors, recycling, regeneration, religions, remotesensing, reports, reprodmed, resources, rheumato, risks, rjpm, robotics, rsee, ruminants, safety, sci, scipharm, sclerosis, seeds, sensors, separations, sexes, shi, signals, sinusitis, siuj, skins, smartcities, sna, societies, socsci, software, soilsystems, solar, solids, spectroscj, sports, standards, stats, std, stratsediment, stresses, surfaces, surgeries, suschem, sustainability, symmetry, synbio, systems, tae, targets, taxonomy, technologies, telecom, test, textiles, thalassrep, therapeutics, thermo, timespace, tomography, tourismhosp, toxics, toxins, tph, transplantology, transportation, traumacare, traumas, tri, tropicalmed, universe, urbansci, uro, vaccines, vehicles, venereology, vetsci, vibration, virtualworlds, viruses, vision, waste, water, welding, wem, wevj, wild, wind, women, world, youth, zoonoticdis

%---------
% article
%---------
% The default type of manuscript is "article", but can be replaced by: 
% abstract, addendum, article, benchmark, book, bookreview, briefcommunication, briefreport, casereport, changes, clinicopathologicalchallenge, comment, commentary, communication, conceptpaper, conferenceproceedings, correction, conferencereport, creative, datadescriptor, discussion, entry, expressionofconcern, extendedabstract, editorial, essay, erratum, fieldguide, hypothesis, interestingimages, letter, meetingreport, monograph, newbookreceived, obituary, opinion, proceedingpaper, projectreport, reply, retraction, review, perspective, protocol, shortnote, studyprotocol, supfile, systematicreview, technicalnote, viewpoint, guidelines, registeredreport, tutorial,  giantsinurology, urologyaroundtheworld
% supfile = supplementary materials

%----------
% submit
%----------
% The class option "submit" will be changed to "accept" by the Editorial Office when the paper is accepted. This will only make changes to the frontpage (e.g., the logo of the journal will get visible), the headings, and the copyright information. Also, line numbering will be removed. Journal info and pagination for accepted papers will also be assigned by the Editorial Office.

%------------------
% moreauthors
%------------------
% If there is only one author the class option oneauthor should be used. Otherwise use the class option moreauthors.

%---------
% pdftex
%---------
% The option pdftex is for use with pdfLaTeX. Remove "pdftex" for (1) compiling with LaTeX & dvi2pdf (if eps figures are used) or for (2) compiling with XeLaTeX.

%=================================================================
% MDPI internal commands - do not modify
\firstpage{1} 
\makeatletter 
\setcounter{page}{\@firstpage} 
\makeatother
\pubvolume{1}
\issuenum{1}
\articlenumber{0}
\pubyear{2026}
\copyrightyear{2026}
%\externaleditor{Firstname Lastname} % More than 1 editor, please add `` and '' before the last editor name
\datereceived{ } 
\daterevised{ } % Comment out if no revised date
\dateaccepted{ } 
\datepublished{ } 
%\datecorrected{} % For corrected papers: "Corrected: XXX" date in the original paper.
%\dateretracted{} % For retracted papers: "Retracted: XXX" date in the original paper.
%\doinum{} % Used for some special journals, like molbank
%\pdfoutput=1 % Uncommented for upload to arXiv.org
%\CorrStatement{yes}  % For updates
%\longauthorlist{yes} % For many authors that exceed the left citation part
%\IsAssociation{yes} % For association journals

%=================================================================
% Add packages and commands here. The following packages are loaded in our class file: fontenc, inputenc, calc, indentfirst, fancyhdr, graphicx, epstopdf, lastpage, ifthen, float, amsmath, amssymb, lineno, setspace, enumitem, mathpazo, booktabs, titlesec, etoolbox, tabto, xcolor, colortbl, soul, multirow, microtype, tikz, totcount, changepage, attrib, upgreek, array, tabularx, pbox, ragged2e, tocloft, marginnote, marginfix, enotez, amsthm, natbib, hyperref, cleveref, scrextend, url, geometry, newfloat, caption, draftwatermark, seqsplit
% cleveref: load \crefname definitions after \begin{document}

%=================================================================
% Please use the following mathematics environments: Theorem, Lemma, Corollary, Proposition, Characterization, Property, Problem, Example, ExamplesandDefinitions, Hypothesis, Remark, Definition, Notation, Assumption
%% For proofs, please use the proof environment (the amsthm package is loaded by the MDPI class).

%=================================================================
% Full title of the paper (Capitalized)
\Title{CADEMAS-ML: A Software Tool for Context-Aware Cooperative and Explainable Automated Decision-Making}

% Author Orchid ID: enter ID or remove command
\newcommand{\orcidauthorA}{0000-0003-3267-6753} % Add \orcidA{} behind the author's name
\newcommand{\orcidauthorB}{0009-0002-0183-7905} % Add \orcidB{} behind the author's name
\newcommand{\orcidauthorC}{0000-0002-7653-1452} % Add \orcidB{} behind the author's name

% Authors, for the paper (add full first names)
\Author{Pavel Novoa-Hernández$^{1,*}$\orcidA{}, Mariia Godz$^{2}$\orcidA{} and David A. Pelta$^{2}$\orcidC{}}

%\longauthorlist{yes}

% MDPI internal command: Authors, for metadata in PDF
\AuthorNames{Pavel Novoa-Hernández, Mariia Godz and David A. Pelta}

%\longauthorlist{yes}

% Affiliations / Addresses (Add [1] after \address if there is only one affiliation.)
\address{%
$^{1}$ \quad Dept. Computer Engineering and Systems, Universidad de La Laguna; pnovoahe@ull.edu.es\\
$^{2}$ \quad Dept. Computer Science and Artificial Intelligence, Universidad de Granada; \{mariiagodz,dpelta\}@ugr.es}

% Contact information of the corresponding author
\corres{Correspondence:pnovoahe@ull.edu.es}

% Current address and/or shared authorship
%\firstnote{Current address: Affiliation.}  
% Current address should not be the same as any items in the Affiliation section.

%\secondnote{These authors contributed equally to this work.}
% The commands \thirdnote{} till \eighthnote{} are available for further notes.

%\simplesumm{} % Simple summary

%\conference{} % An extended version of a conference paper

% Abstract (Do not insert blank lines, i.e. \\) 
\abstract{Automated decision-making systems are increasingly used to support prediction and classification tasks. In practice, however, decisions rarely depend on predictive models alone. Organizations must also consider contextual factors such as strategic priorities, resource availability, operational constraints, and policy objectives, many of which cannot be captured directly from data. While machine learning provides predictive power and fuzzy systems offer a flexible way to represent contextual knowledge, existing solutions provide limited support for integrating both perspectives within a transparent and reusable decision-making process.
This paper presents CADEMAS-ML, an open-source software tool that operationalizes the Cooperative Automated Decision-Making Systems (CADEMAS) framework for context-aware case prioritization. The tool enables multiple predictive models to cooperate in the assessment of cases while incorporating contextual information through configurable fuzzy knowledge bases. In addition to combining prediction and context, CADEMAS-ML provides mechanisms to evaluate the robustness of prioritization results under different decision policies and to explain how predictive and contextual factors contribute to each recommendation. These capabilities support transparency, auditability, and informed decision-making in settings where priorities may change across organizational scenarios.
The tool is validated through an employee attrition case study involving multiple predictive models and contrasting organizational contexts. Consistent with a previous instantiation of the CADEMAS framework, the results show that prioritization outcomes are influenced not only by predicted risk levels but also by the contextual conditions under which decisions are made. They further demonstrate how the proposed approach enables decision-makers to understand, justify, and analyze ranking outcomes before actions are taken. Overall, CADEMAS-ML provides a practical and reusable environment for designing, evaluating, and explaining cooperative and context-aware automated decision-making systems.}
% Keywords
\keyword{automated decision-making systems; fuzzy decision making; context-aware decision support; machine learning; prioritization} 

%%%%%%%%%%%%%%%%%%%%%%%%%%%%%%%%%%%%%%%%%%
\begin{document}

\section{Introduction}\label{sec:intro}

Automated decision-making (ADM) systems have become an integral part of modern organizations. Following the broad understanding of ADM as computational processes that produce ratings, recommendations, decisions, predictions, or related outputs from data and predefined objectives \citep{EIP2022}, such systems may include machine-learning classifiers, rule-based procedures, optimization models, simulation components, or other algorithmic decision units. In practice, predictive and classification models are among the most common instantiations, supporting tasks such as credit scoring, fraud detection, medical triage, customer churn analysis, and human resources management.

Despite these advances, an important limitation remains. The outputs generated by ADM units, such as recommendations, ratings, class probabilities, or risk scores, are often interpreted in isolation from the organizational environment in which decisions are ultimately made \citep{Tariq2024a,Deliu2026}. In practice, a technically plausible automated output rarely provides sufficient guidance for action. The appropriate response depends on contextual factors such as organizational priorities, resource constraints, departmental objectives, regulatory constraints, and prevailing strategic conditions. As a result, identical automated assessments may lead to different decisions under different circumstances.

This observation highlights the cooperative nature of real-world decision-making \citep{Tariq2024a}. Decisions are seldom based on a single decision agent or automated component. Organizations frequently rely on heterogeneous units developed for different departments, objectives, data sources, or stakeholder groups. Their outputs must therefore be harmonized into a coherent assessment before they can support action. At the same time, domain experts contribute contextual knowledge that cannot always be represented through precise numerical rules. Such knowledge is often expressed through gradual concepts and linguistic assessments that reflect uncertainty, preference, and strategic intent.

Fuzzy logic provides a well-established mathematical framework for representing this type of contextual reasoning \citep{PerezCanedo2023,Lamata2020}. Through fuzzy sets and linguistic variables, experts can model concepts such as urgency, strategic relevance, or suitability in a transparent and interpretable manner. This makes fuzzy approaches particularly attractive for decision-support environments in which quantitative predictions must be complemented by qualitative contextual considerations.

To address these challenges, \citet{NovoaHernandez2026} proposed \textit{CADEMAS} (Cooperative Automated Decision-Making Systems), a conceptual framework for coordinating heterogeneous automated decision units. The framework separates the generation of decision evidence from its contextual evaluation, allowing outputs from multiple decision units to be integrated and subsequently assessed according to scenario-specific conditions and priorities.

The practical value of this approach was subsequently demonstrated by \citet{Godz2026}, who instantiated CADEMAS in the context of employee-attrition prevention. In that work, predictive machine-learning models developed from different organizational perspectives were treated as cooperating decision units and combined with contextual organizational policies to support case prioritization. The results showed that the framework could effectively integrate predictive evidence and contextual knowledge within a coherent decision-making process.

Although CADEMAS provides a solid conceptual foundation for cooperative and context-aware decision-making, its practical application typically involves a non-trivial implementation effort. In particular, users must define how case data are distributed across multiple decision units, how their outputs are collected and aligned, how these outputs are integrated, and how the resulting decision proposal is evaluated under explicit contextual assumptions. While these steps are conceptually well-defined, they must be repeatedly specified and coordinated when the framework is applied to new domains or datasets, which introduces operational overhead in practice for both researchers and decision-makers.

Beyond this implementation effort, decision-support settings also require additional capabilities that are often critical in practice, such as assessing the robustness of decisions under alternative policy configurations and providing clear explanations of how individual recommendations are derived. These aspects are particularly relevant in organizational environments where decisions must be justified, audited, and analyzed under changing contextual conditions. However, these functionalities are not explicitly addressed in the original CADEMAS formulation nor in subsequent application-focused instantiations.

This paper addresses these requirements by presenting \textbf{CADEMAS-ML}, a software tool that operationalizes the CADEMAS framework for cooperative and context-aware automated decision-making. The main contributions of this work are as follows:
\begin{itemize}
\item The development of CADEMAS-ML, an open-source software tool that implements the CADEMAS framework for context-aware cooperative decision-making using machine-learning models as automated decision units.
\item The integration of predictive analytics, fuzzy contextual reasoning, and explainable decision support within a unified prioritization process, enabling transparent justification and case-level interpretability.
\item A robustness-analysis module that allows decision-makers to assess the sensitivity and stability of prioritization outcomes under alternative contextual policies and preference configurations.
\item An empirical validation in an employee-attrition prevention scenario involving multiple predictive models and organizational contexts.
\end{itemize}

The explainability and robustness capabilities introduced in this work are central to its operational contribution. They complement the original CADEMAS formulation, which is primarily architectural, as well as prior application-oriented instantiations, by supporting an interactive decision-support workflow in which users can inspect the rationale behind individual prioritizations and assess their stability under alternative contextual assumptions. This shifts the focus from producing a single prioritization outcome to enabling its interpretation, justification, and analysis in practice.

The remainder of the paper is organized as follows. Section~\ref{sec:background} reviews the CADEMAS framework and positions it within related cooperative and context-aware decision-support approaches. Section~\ref{sec:tool} presents CADEMAS-ML, detailing its architecture and functionalities and illustrating its use in an employee-attrition case study. Finally, Section~\ref{sec:conclusions} concludes the paper and outlines directions for future work.

\section{Background and related works}\label{sec:background}

This section sets out the conceptual foundations on which CADEMAS-ML builds. We first revisit the CADEMAS framework and formalize it for the prioritization setting addressed here (Section~\ref{sec:cademas}). We then situate the tool among related cooperative-decision frameworks and application-oriented decision-support systems, clarifying the gap it is designed to fill (Section~\ref{sec:related}).

\subsection{CADEMAS framework}\label{sec:cademas} 

As introduced in Section~\ref{sec:intro}, CADEMAS \citep{NovoaHernandez2026} provides a general structure for designing systems in which multiple automated decision-making units cooperate to produce contextually aligned decisions. As noted there, its motivation is not limited to prediction: organizational decisions rarely emerge from a single autonomous unit with a single representation of the problem, and may instead depend on heterogeneous signals produced by classifiers, optimization models, simulation components, rule-based systems, large-language-model agents, or institutional procedures, each with its own assumptions, data access, output type, and reliability level.

To formalize this idea for the prioritization setting addressed in this paper, let $\mathcal{X} = \{x_1, \dots, x_n\}$ denote the case dataset, and let $\mathcal{A} = \{A_1, \dots, A_K\}$ represent a collection of ADM units. In the general CADEMAS framework, each $A_k$ may produce an output $o_k(x_i)$ whose semantics are not necessarily probabilistic: it may be a class, a score, a ranking, a symbolic recommendation, a rule activation, or another decision object. A decision-integration layer is therefore responsible for transforming these heterogeneous outputs into a shared representation and synthesizing them into an integrated proposal. In the ML specialization implemented by CADEMAS-ML, the ADM units are predictive models and the harmonized output is a scalar risk or priority signal $p_k(x_i) \in [0,1]$. Since these components may differ in performance and trustworthiness, each $A_k$ is assigned a confidence weight $w_k \geq 0$, with $\sum_{k=1}^{K} w_k = 1$. These weights can be derived from validation metrics, expert judgment, or governance constraints. The ML specialization aggregates individual outputs into an integrated risk score:
\begin{equation}
R(x_i) = \sum_{k=1}^{K} w_k \, p_k(x_i).
\end{equation}

Rather than treating the integrated output as the final decision criterion, CADEMAS introduces an explicit contextual layer that captures organizational priorities, policies, constraints, and scenario-dependent considerations. The original framework does not prescribe a single context formalism. Context may be represented through crisp rules, utility functions, regulatory constraints, fuzzy propositions, or other evaluative mechanisms. In CADEMAS-ML, and following the attrition application of \citet{Godz2026}, fuzzy logic is used because it supports gradual degrees of satisfaction instead of binary rules. Let $\mathcal{C} = \{C_1, \dots, C_L\}$ be a set of contextual factors. For each case $x_i$, the degree to which it satisfies factor $C_\ell$ is given by $\mu_\ell(x_i) \in [0,1]$, and these values are aggregated into a contextual score:
\begin{equation}
C(x_i) = H\big(\mu_1(x_i), \dots, \mu_L(x_i)\big),    
\end{equation}
where $H: [0,1]^L \rightarrow [0,1]$ is an aggregation operator that can take different forms depending on the decision policy (e.g., t-norms, averages, or hierarchical compositions).

The final prioritization index in CADEMAS-ML is obtained by combining the ML-derived risk and the contextual suitability score:
\begin{equation}
\pi(x_i) = \Psi\big(R(x_i), C(x_i)\big),    
\end{equation}
where $\Psi$ is a monotonic function that defines the trade-off between both components. A key design principle inherited from CADEMAS is the separation between the integrated ADM proposal and its contextual interpretation. In the ML specialization, this makes explicit what the predictive models suggest and what the organization considers relevant.

As mentioned before, the framework has been instantiated in employee-attrition prevention by \citet{Godz2026}. In that application, multiple machine-learning models were trained on different organizational perspectives (e.g., Human Resources, Finance, Operations, and Employee Wellbeing). Each model acts as an ADM unit within $\mathcal{A}$, and their outputs are combined using weights derived from validation performance, specifically AUC. In parallel, the contextual layer encodes scenario-dependent organizational priorities (e.g., economic crisis or digital transformation settings) using fuzzy membership functions and rule-based structures.

Although this instantiation demonstrates the practicality of the framework, it represents only one possible configuration. In general, CADEMAS is designed to remain flexible across several dimensions, including the type of automated decision units, the semantics of their outputs, the mechanism used for aligning heterogeneous decision signals, the integration strategy, and the way contextual information is incorporated into the final decision.

As will be detailed later, CADEMAS-ML implements specific design choices for these components in order to provide a concrete and reproducible software realization of the framework. Within the general CADEMAS formulation, however, these elements remain configurable: confidence weights can be derived from different performance criteria, aggregation strategies may reflect alternative decision policies, and the final prioritization can be adjusted to capture different trade-offs between predictive evidence and contextual constraints.

This configurability is essential for adapting CADEMAS to heterogeneous organizational environments and motivates the need for a software layer that operationalizes these design decisions in a systematic and reusable manner.

\subsection{Related frameworks and decision-support applications}\label{sec:related}

Recent research has increasingly recognized that effective decision support requires more than accurate isolated models. In many application domains, decisions emerge from the interaction between automated units, expert knowledge, and contextual considerations that may change over time. For clarity, it is useful to distinguish two related but different bodies of work: general cooperative-decision frameworks, which address how multiple decision entities interact, and application-oriented decision-support systems, which implement concrete methods for a domain or family of decision problems.

At the framework level, \citet{Tariq2024a} propose an AI-to-collaborator (A2C) framework in which humans and AI systems share decision authority through explicit interaction mechanisms. Related work on hybrid intelligence emphasizes that authority distribution and human interpretive capacity should be modeled separately, because the level of AI integration and the cognitive flexibility of human decision-makers jointly shape decision quality and accountability \citep{Deliu2026}. A similar pattern appears in strategic management, where AI contributes large-scale analytical capability but managers retain interpretive authority over final choices \citep{Telaumbanua2025}. A neighboring line of work studies cooperation among multiple AI agents, including LLM-based multi-agent reinforcement learning \citep{Sun2024,Liu2025} and trust-weighted consensus mechanisms \citep{Shakya2026}. These contributions are important alternatives or neighboring frameworks because they address cooperation among decision entities. However, their main focus is not the reusable implementation of a CADEMAS-like workflow in which heterogeneous ADM outputs are explicitly separated from context modulation and final prioritization.

At the application and decision-support level, substantial progress has been made in fuzzy decision support, multi-criteria decision making, and hybrid fuzzy-ML workflows. Applications such as \textsc{MakeDecision} \citep{Wieckowski2024} and the fuzzy DEMATEL-based tool of \citet{Chekry2024} emphasize transparency and traceability by exposing intermediate calculations to users. A related stream has explored hybrid pipelines that combine expert judgment, causal weighting, and predictive components in domain-specific settings \citep{Abdulgader2018,Sahraei2022,Ozkurt2025,Mikhaylov2026}. More recent work has also extended the fuzzy toolbox with richer aggregation operators for elicited expert preferences \citep{Suvitha2025}. These contributions demonstrate the value of combining data-driven evidence with contextual or expert-defined criteria, but most are tailored to particular methods or domains rather than organized around the CADEMAS decomposition of ADM units, integration, context, and final prioritization.

Open-source fuzzy and MCDM libraries, including \textsc{pyFDM} \citep{Wieckowski2022}, \textsc{PyIFDM} \citep{Wieckowski2023}, \texttt{FuzzyR} \citep{Chen2020}, the \textsc{IFMCDM} package of \citet{Roszkowska2024}, the \textsc{Eligere} platform \citep{Grazioso2023}, and the intuitionistic fuzzy dimensional-analysis approach of \citet{PerezDominguez2024}, are also relevant. They improve reproducibility and provide reusable computational building blocks. 

%Nevertheless, software availability is not the central distinction pursued here. The key issue is whether a system captures the cooperative ADM problem that CADEMAS is intended to solve: multiple decision-producing units, explicit integration of their outputs, contextual evaluation that is kept conceptually separate, and a final decision mechanism whose assumptions can be inspected and varied.

\begin{table}[t]
\caption{Comparison of representative application-oriented decision-support works related to CADEMAS-ML.\label{tab:related}}
\small
\begin{adjustwidth}{-\extralength}{0cm}
	\begin{tabularx}{\fulllength}{>{\raggedright\arraybackslash}p{2cm} >{\raggedright\arraybackslash}X p{2cm} p{2cm} >{\raggedright\arraybackslash}X}
		\toprule
		\textbf{Work} & \textbf{Decision-support focus} & \textbf{Multiple ADM units} & \textbf{Explicit context layer} & \textbf{Relation to CADEMAS-ML} \\
		\midrule
		\citet{Godz2026} & Employee-attrition prevention using CADEMAS with departmental ML models & Yes & Yes & Direct application that motivates the implemented workflow \\
		\citet{Tariq2024a} & Human--AI collaborative decision framework with automated, augmented, and collaborative modes & Partly & Partly & Strong framework-level precedent, but not a configurable fuzzy/ML prioritization workflow \\
		\citet{Wieckowski2024} & Auditable fuzzy MCDM through a web-based tool & No & Partly & Strong transparency, but no cooperative ADM integration layer \\
		\citet{Chekry2024} & Open-source fuzzy DEMATEL decision-support tool & No & Partly & Traceable fuzzy causal analysis, but method-specific rather than CADEMAS-like \\
		\citet{Ozkurt2025} & Hybrid fuzzy MCDM, ML, and XAI for transparent ranking and prediction & Partly & Partly & Combines ranking and prediction, but lacks a general cooperative ADM architecture \\
		\citet{Sahraei2022} & Fuzzy DEMATEL/BWM/ANP pipeline for spatial risk prioritization & No & Partly & Explicit causal weighting under uncertainty, but application-specific \\
		\midrule
		CADEMAS-ML (this work) & Configurable prioritization with ML-based ADM units and fuzzy context scenarios & Yes & Yes & Operationalizes CADEMAS for reusable ML-driven decision support \\
		\bottomrule
	\end{tabularx}
\end{adjustwidth}
\end{table}

Table~\ref{tab:related} summarizes representative application-oriented contributions in these areas. Two observations emerge from this comparison. First, published systems already provide strong ingredients for transparent decision support, including deferral-aware collaboration, auditable fuzzy MCDM, hybrid ranking pipelines, and reusable software libraries \citep{Tariq2024a,Wieckowski2024,Ozkurt2025,Wieckowski2022}. Second, these ingredients are usually developed in isolation. What remains missing is an operational workflow in which multiple ADM units are treated as cooperating components, their outputs are integrated explicitly, contextual evaluation is modeled as a separate layer, and the final prioritization can be recomputed under different decision policies.

CADEMAS-ML is designed to address this gap for the predictive-model setting. Building on the framework and its attrition instantiation discussed above, it provides a configurable software implementation in which ML-based ADM outputs and contextual reasoning coexist as explicit and auditable components of the decision-making process. Its planned open-source availability is an additional benefit, but the main contribution is the operationalization of the CADEMAS decomposition in a reusable decision-support tool.


\section{CADEMAS-ML software tool}\label{sec:tool}

This section presents CADEMAS-ML as a software realization of one concrete CADEMAS specialization: cooperative prioritization supported by multiple predictive ADM units and an explicit fuzzy context layer. We first explain how the general CADEMAS architecture is instantiated in a configurable ML workflow. Special attention is given to the fuzzy context model, the case-level \textit{Explain} module, and the robustness analysis of the prioritization ranking. We then use the employee-attrition example distributed with the repository to show the main functionalities of the application.

\subsection{Architecture and improvements}\label{sec:architecture_improvements}

CADEMAS-ML is not a domain-specific attrition script. It is an operational implementation of CADEMAS for the case in which the cooperating ADM units are predictive models. In line with \citet{NovoaHernandez2026}, the software keeps partial automated assessments separate from their contextual interpretation. Its contribution is to make this separation executable for ML-based decision support. Predictive components, contextual rules, and decision policies are configured explicitly and can be inspected through the interface. Thus, the attrition study of \citet{Godz2026} becomes one possible configuration of the tool, not its only intended use.

\begin{figure}[t]
\centering
\includegraphics[width=\textwidth]{figures/architecture.pdf}
\caption{Operational architecture of CADEMAS-ML application. The software instantiates the CADEMAS \citep{NovoaHernandez2026} architecture as a configurable vertical pipeline in which predictive risk and contextual alignment are computed independently before being integrated into a prioritization index.}
\label{fig:cademas_arch}
\end{figure}

The application is implemented in Python and combines a Streamlit \citep{streamlit} front end, H2O\citep{H2O.ai2026} inference, and a vectorized fuzzy-context engine. Interactive visual analytics are rendered with Altair \citep{VanderPlas2018} and embedded in the \textit{Overview}, \textit{Context}, \textit{Explain}, and \textit{Robustness} modules. 

%Because Altair produces declarative chart specifications, the resulting figures---including scatter plots, membership-function views, attribution bar charts, and the three robustness summaries introduced below---can be exported directly for reporting and publication. 

A run starts from four inputs: the case dataset, the MOJO models, a model-metadata JSON file, and one or more context JSON files. These inputs map naturally onto the CADEMAS layers. The dataset and configuration files act as the input manager. The MOJO files instantiate the ADM units in this ML specialization. Their weighted probabilistic outputs form the integrated predictive risk. The fuzzy-rule engine provides the contextual model. The final layer combines both streams through a user-adjustable policy parameter. Figure~\ref{fig:cademas_arch} summarizes this architecture.



\subsubsection{ML-based cooperative risk}

The predictive layer follows the cooperative aggregation principle of CADEMAS as specialized in the attrition instantiation of \citet{Godz2026}. CADEMAS-ML exposes that specialization as a reusable software component. Let $\mathcal{X}=\{x_1,\ldots,x_n\}$ be the case dataset and let $\mathcal{M}=\{M_1,\ldots,M_K\}$ be the uploaded predictive ADM units. Each model $M_k$ produces, for case $x_i$, a positive-class probability
\[
\hat{p}_{ik} \in [0,1].
\]
The models need not use the same feature subset. Each MOJO receives the available case matrix and uses its internal schema to select the required variables. The contribution of each component is controlled by a non-negative confidence weight $w_k$. This weight is computed from a user-selected validation metric $s_k$:
\begin{equation}
w_k = \frac{s_k}{\sum_{\ell=1}^{K} s_\ell}, \qquad \sum_{k=1}^{K}w_k=1.
\label{eq:model_weights}
\end{equation}
The cooperative machine-learning risk score is then
\begin{equation}
R_i=\sum_{k=1}^{K}w_k\hat{p}_{ik}.
\label{eq:ml_risk}
\end{equation}
This mathematical form is the ML-specific realization of the more general CADEMAS decision-integration layer. The improvement introduced by CADEMAS-ML is practical configurability. Users can upload any number of compatible MOJO models, read model metadata from JSON, and choose the validation metric used for weighting.

\subsubsection{Configurable fuzzy context model}

The context layer is where CADEMAS-ML most clearly turns the conceptual framework into an operational artifact. The role of context follows \citet{NovoaHernandez2026}, and the use of fuzzy propositions follows the attrition application in \citet{Godz2026}. The difference is that CADEMAS-ML does not hard-code those contexts. It represents them as JSON specifications. The same formalism can therefore be reused with other domains, rule families, and decision policies.

Formally, a context is defined as a collection of fuzzy propositions. Let $V_r$ be a case attribute and let $\tilde{F}_r$ be a fuzzy set representing a linguistic predicate over the universe of $V_r$. For case $x_i$, an atomic rule $r$ evaluates the proposition ``$V_r$ is $\tilde{F}_r$'' as
\begin{equation}
\mu_r(x_i)=\mu_{\tilde{F}_r}(x_{ir}) \in [0,1],
\label{eq:atomic_mu}
\end{equation}
where $x_{ir}$ denotes the value of attribute $V_r$ for case $x_i$. The implementation supports triangular, trapezoidal, linear increasing, and linear decreasing membership functions. It also supports categorical maps and categorical sets. The former assign degrees directly to symbolic values. The latter group categories into linguistic labels that can be scored later in derived rules.

Atomic rules may include activation conditions. If $m_r(x_i)\in\{0,1\}$ denotes whether the \texttt{when} condition of rule $r$ is satisfied for case $x_i$, the effective membership is
\begin{equation}
\mu_r^{*}(x_i)=m_r(x_i)\mu_r(x_i),
\label{eq:conditional_mu}
\end{equation}
with $m_r(x_i)=1$ for unconditional rules. This makes it possible to represent policies such as department-specific salary suitability without duplicating datasets or writing procedural code.

Beyond atomic rules, CADEMAS-ML allows the definition of derived rules. A derived rule $d$ combines a subset of already computed memberships through an aggregation operator $G_d$:
\begin{equation}
\mu_d(x_i)=G_d\left(\mu_{r_1}^{*}(x_i),\ldots,\mu_{r_q}^{*}(x_i)\right).
\label{eq:derived_mu}
\end{equation}
The implemented operators are minimum-based \texttt{AND}, maximum-based \texttt{OR}, algebraic \texttt{PRODUCT}, and arithmetic \texttt{AVERAGE}. The context alignment score is then obtained by applying a final aggregation policy $H$ to selected atomic and derived memberships:
\begin{equation}
C_i = H\left(\mu_1(x_i),\ldots,\mu_L(x_i)\right).
\label{eq:context_score}
\end{equation}
The function $H$ can be selected interactively or specified as a declarative logic tree in the context JSON. This captures the attrition instantiation in \citet{Godz2026}, but makes it portable. It also lets organizations express conservative, compensatory, hierarchical, or conditional contextual policies without changing the application code. The closed-form specification of the context used in the case study is reported in Appendix~\ref{app:context_specs}.

All atomic and derived memberships are retained as an audit matrix. The final context score is therefore not a black-box modifier of the ML risk. It can be decomposed into the rule activations that produced it.

\subsubsection{Decision integration and robustness}

Predictive evidence and contextual suitability remain separate until the final decision stage. For a fixed policy parameter $\lambda\in[0,1]$, CADEMAS-ML computes the prioritization score as
\begin{equation}
P_i(\lambda)=\lambda R_i+(1-\lambda)C_i.
\label{eq:priority_lambda}
\end{equation}
The interpretation is direct. A value $\lambda=1$ produces a risk-driven ranking. A value $\lambda=0$ produces a context-driven ranking. Intermediate values encode explicit trade-offs between both components.

The robustness module is a specific improvement of CADEMAS-ML over the previous static attrition instantiation. Let $\Lambda=\{\lambda_0,\ldots,\lambda_q\}$ be a finite grid of policy values. For each $\lambda\in\Lambda$, the tool computes the score vector $P(\lambda)$ and the ranking
\begin{equation}
r_i(\lambda)=1+\left|\{j: P_j(\lambda)>P_i(\lambda)\}\right|,
\label{eq:rank_lambda}
\end{equation}
with deterministic tie-breaking by case order. This produces a rank trajectory for every case. CADEMAS-ML summarizes these trajectories through three complementary Altair views: a bump chart, a ranks box plot, and a rank-acceptability heatmap. The bump chart shows how each case moves as $\lambda$ changes. The box plot summarizes the distribution of ranks over $\Lambda$; cases are ordered by median rank and interquartile range. Finally, the rank acceptability index estimates how often a case attains a given rank:
\begin{equation}
b_i^r = \frac{1}{|\Lambda|}\sum_{\lambda\in\Lambda}\mathbf{1}\{r_i(\lambda)=r\}.
\label{eq:rank_acceptability}
\end{equation}
Together, these indicators distinguish stable priorities from policy-sensitive cases.

\subsubsection{Case-level explanation and natural-language recommendations}

A second improvement introduced by CADEMAS-ML is the \textit{Explain} module. The original CADEMAS framework specifies the separation between ADM outputs, context evaluation, and final decision integration, but it does not prescribe a case-level explanation mechanism. Similarly, the attrition instantiation of \citet{Godz2026} reports contextual prioritization outcomes, but does not provide an interactive module that explains individual cases in natural language. CADEMAS-ML fills this gap by storing intermediate artifacts during the analysis run and exposing them through a four-part local explanation.

For a selected case $x_i$, the first explanation level decomposes the final score in Equation~\eqref{eq:priority_lambda} into its two additive components:
\begin{equation}
P_i(\lambda)=\underbrace{\lambda R_i}_{\text{predictive contribution}}+
\underbrace{(1-\lambda)C_i}_{\text{contextual contribution}}.
\label{eq:priority_decomposition}
\end{equation}
This decomposition makes explicit whether the prioritization is mainly driven by predictive risk, contextual alignment, or a balance between both.

The second level provides local ML attribution. Since the uploaded H2O AutoML MOJO models are treated as black-box predictive ADM units and may include stacked ensembles, CADEMAS-ML uses a model-agnostic perturbation procedure rather than relying on model-specific SHAP values. For each model $M_k$ and feature $f$, the tool compares the original predicted probability with the prediction obtained after replacing the value of $f$ by a cohort baseline $b_f$ (median for numerical variables and mode for categorical variables):
\begin{equation}
\Delta_{ikf}=\hat{p}_{ik}(x_i)-\hat{p}_{ik}(x_i^{(f \leftarrow b_f)}).
\label{eq:local_attribution_model}
\end{equation}
The model-specific effects are aggregated with the same confidence weights used for the cooperative risk:
\begin{equation}
\Delta_{if}=\sum_{k=1}^{K} w_k \Delta_{ikf}.
\label{eq:local_attribution_global}
\end{equation}
Positive values indicate features whose observed value increases the local predicted risk relative to the cohort baseline, while negative values indicate features whose observed value decreases it. These quantities are expressed in probability points and should be interpreted as approximate local effects, not as causal impacts.

The third level traces the contextual component by presenting the fuzzy memberships $\mu_r(x_i)$ of the atomic and derived rules that produced $C_i$. This connects the numerical context score with the policy rules encoded in the JSON context specification. Finally, the fourth level generates a deterministic natural-language summary. The text combines the priority score, the balance parameter $\lambda$, the risk and context levels, the most relevant ML drivers, the most relevant contextual drivers or bottlenecks, and a generic tier-based recommendation. The recommendation uses the final score $P_i$: \textit{Tier 1} for high priority, \textit{Tier 2} for medium or borderline priority requiring manual qualitative review, and \textit{Tier 3} for low priority under the current prediction and context guidelines.

\subsection{Main functionalities}\label{sec:case}\label{sec:casestudy}

We illustrate the main functionalities of CADEMAS-ML through an employee-attrition scenario that reproduces the methodological setup of \citet{Godz2026} as an interactive software workflow. The objective is not to improve attrition prediction accuracy, but to show how cooperative predictive assessments and explicit contextual policies are configured, executed, audited, and explained within the architecture described in Section~\ref{sec:architecture_improvements}.

The scenario assumes that an organization must prioritize retention interventions among employees at attrition risk when resources are limited. Predictive evidence is not monolithic: we simulate four organizational perspectives---Culture and Wellbeing, Finance, Operations and Development, and Human Resources---each modeled as an independent ADM unit that observes only a subset of employee attributes from the same workforce records. This hypothetically reflects departments with partial, perspective-specific access to human-resources data. All units were trained offline on the publicly available Kaggle dataset \textit{IBM HR Analytics Employee Attrition and Performance}\footnote{\url{https://www.kaggle.com/datasets/pavansubhasht/ibm-hr-analytics-attrition-dataset} (accessed 19 March 2026).}, which contains 1,470 employee records and 35 attributes describing demographic, professional, compensation, performance, and organizational conditions. Following \citet{Godz2026}, H2O AutoML was run independently for each feature subset defined in Table~\ref{tab:attrition_feature_assignment}, using an 80/20 train-validation partition and a one-hour time budget per model. The selected MOJO artifacts, their validation metrics, and feature assignments are distributed with the application. At run time, CADEMAS-ML loads these models, produces case-level probabilities, and aggregates them into the cooperative risk score $R_i$ using the AUC-based weights in Equation~\eqref{eq:model_weights}.

Contextual prioritization is encoded separately through JSON fuzzy specifications. The example bundle includes two alternative organizational policies: \textit{Economic Crisis} and \textit{Digital / Technological Transformation}. They share the same predictive layer but differ in the rules used to compute $C_i$. In the walkthrough below, we focus on \textit{Digital / Technological Transformation}, which prioritizes upskilling capacity, role adaptability, career scalability, and investment horizon during a technological-change program. Appendix~\ref{app:context_specs} reports the corresponding specification in closed mathematical form.

For interactive demonstration, the repository provides a 20-case CSV subset with the same 35 columns as the original dataset, including the \texttt{attrition} outcome label. This label is retained for verification but is not used as a prioritization criterion. Employee names in \texttt{Case\_ID} were randomly generated to simplify reference in the interface and in this paper. The subsections that follow describe how CADEMAS-ML processes these inputs through its interface modules: loading configuration, computing $R_i$ and $C_i$, integrating them into $P_i(\lambda)$, and supporting inspection through the \textit{Overview}, \textit{Models}, \textit{Context}, \textit{Explain}, and \textit{Robustness} views.

\subsubsection{Input configuration and analysis workflow}

The entry point is the \textit{Home} module and the sidebar upload controls (Figure~\ref{fig:home_upload}). The user provides the four inputs required by the architecture---model-metadata JSON, context JSON, MOJO models, and the case dataset---or loads the ready-made example distributed under \texttt{example\_attrition/}\footnote{\url{https://github.com/pnovoa/cademas-app/tree/main/example_attrition}.}. The application \textit{Help} panel documents the same bundle. Before execution, the user selects the validation metric used for model weighting, chooses one of the uploaded context specifications, sets the balance parameter $\lambda$, and triggers \textit{Run analysis}. CADEMAS-ML then executes the pipeline defined in Section~\ref{sec:architecture_improvements}: MOJO inference, fuzzy context evaluation, and score integration. Table~\ref{tab:tool_model_weights} summarizes the four pre-trained ADM units included in the example and the AUC-based weights recomputed during the run.

\begin{figure}[t]
\centering
\includegraphics[width=\textwidth]{figures/home.png} 
\caption{\textit{Home} and input configuration module. The user uploads or loads the four inputs and launches the analysis pipeline described in Section~\ref{sec:architecture_improvements}.}
\label{fig:home_upload}
\end{figure}

Following the training protocol of \citet{Godz2026}, the four MOJO models correspond to the organizational perspectives summarized in Table~\ref{tab:tool_model_weights}. Using AUC as the weighting metric, the reproduced run assigns a lower weight to Operations and Development because its validation AUC is lower than that of the other three components.

\begin{table}[t]
\caption{Models used in the attrition example and AUC-based weights computed by CADEMAS-ML. The \textit{AutoML model} column summarizes the artifact exported by H2O AutoML.\label{tab:tool_model_weights}}
\small 
\begin{tabularx}{\textwidth}{lllcc}
\toprule
\textbf{ID} & \textbf{Organizational perspective} & \textbf{AutoML model} & \textbf{AUC} & \textbf{Weight} \\
\midrule
$M_1$ & Culture and Wellbeing & Stacked Ensemble (BestOfFamily) & 0.9941 & 0.2623 \\
$M_2$ & Finance & Stacked Ensemble (BestOfFamily) & 0.9972 & 0.2631 \\
$M_3$ & Operations and Development & GBM (grid 1, model 26) & 0.8022 & 0.2116 \\
$M_4$ & Human Resources & Stacked Ensemble (BestOfFamily) & 0.9968 & 0.2630 \\
\bottomrule
\end{tabularx}
\end{table}

\subsubsection{\textit{Overview}: prioritized cases}

After execution, the \textit{Overview} module surfaces the outputs of Equations~\eqref{eq:ml_risk}--\eqref{eq:priority_lambda}. Figure~\ref{fig:overview} shows the exported interface for the \textit{Digital / Technological Transformation} run at $\lambda=0.5$. At the top, the \textit{Prioritization: Risk vs Context} scatter plot locates each case in the $(R_i,C_i)$ plane and encodes $P_i$ through colour; the adjacent \textit{Prioritization Score} histogram summarizes the cohort-wide distribution of prioritization scores. Below, the \textit{Prioritized Cases} table lists all cases ordered by $P_i$, reporting $R_i$, $C_i$, and the \texttt{attrition} outcome label retained for verification only. In this run, $C_i$ ranged from 0.1500 to 0.7100, with a mean of 0.4663. The ten highest-priority cases at $\lambda=0.5$ are used later in the robustness analysis of Figure~\ref{fig:robustness}.

\begin{figure}[H]
\centering
\includegraphics[width=\textwidth]{figures/overview.png}
\caption{\textit{Overview} module for the attrition example at $\lambda=0.5$ in the \textit{Digital / Technological Transformation} context. The scatter plot and histogram summarize the risk--context landscape and the distribution of $P_i$; the \textit{Prioritized Cases} table below lists the full cohort by prioritization score, including high-priority employees and context-aligned but low-risk cases such as \textit{Noah Lewis}.}
\label{fig:overview}
\end{figure}

The prioritization ranking in Fig.~\ref{fig:overview} is driven primarily by cooperative predictive risk. The highest-priority cases exhibit high $R_i$, while $C_i$ varies more widely. \textit{Evelyn Taylor} leads the ranking because both components are high ($R_i=0.9555$, $C_i=0.7000$). Cases such as \textit{Lucas Wright} and \textit{Abigail Anderson} remain near the top mainly because of critical predictive risk, even though their contextual alignment is moderate or low. By contrast, \textit{Noah Lewis} appears far down the table despite having the highest $C_i$ in the cohort ($C_i=0.7100$): his cooperative risk is low ($R_i=0.0651$), yielding $P_i(0.5)=0.3875$. The juxtaposition of scatter plot, score distribution, and full case table illustrates why CADEMAS-ML reports $P_i$, $R_i$, and $C_i$ separately rather than ranking cases by contextual alignment alone.

\subsubsection{\textit{Models}: inspecting cooperative risk}

The \textit{Models} module exposes the intermediate results of the cooperative risk layer defined in Equation~\eqref{eq:ml_risk}. Figure~\ref{fig:models} shows the exported interface for the reproduced run. The upper table, \textit{Model Weights}, links each uploaded MOJO to its normalized contribution $w_k$; the values displayed are consistent with those reported in Table~\ref{tab:tool_model_weights}. The lower table, \textit{Risk Probabilities}, reports the per-case positive-class probability returned by each model and the cooperative score $R_i$ in the column immediately after \texttt{CaseID}; cases are ordered by decreasing $R_i$.

\begin{figure}[t]
\centering
\includegraphics[width=\textwidth]{figures/models.png}
\caption{\textit{Models} module for the attrition example, showing the model-weight table that assigns each MOJO its contribution $w_k$ and, below, the risk-probability table with Global ML Risk ($R_i$), per-model probabilities, and cases sorted by decreasing $R_i$.}
\label{fig:models}
\end{figure}

In the exported view, \textit{Lucas Wright} appears first ($R_i=0.9882$), followed by \textit{Amelia Davis} ($R_i=0.9879$) and \textit{Theodore Adams} ($R_i=0.9803$). This ordering reflects the dominance of cooperative predictive risk in the cohort and differs from the prioritization ranking in Fig.~\ref{fig:overview}, where \textit{Evelyn Taylor} leads because $P_i$ also weights contextual alignment. The side-by-side per-model columns make the cooperative aggregation auditable: although all four MOJOs assign high attrition probabilities to the top rows, Operations and Development contributes less to $R_i$ because its validation weight is lower ($w_3=0.2116$, consistent with Table~\ref{tab:tool_model_weights}). Together, the two tables connect the model-confidence layer of Equation~\eqref{eq:model_weights} to the case-level cooperative risk fed into the prioritization score.

\subsubsection{\textit{Context}: auditing fuzzy policies}

The \textit{Context} module renders the fuzzy audit underlying Equation~\eqref{eq:context_score}. Figure~\ref{fig:context} exports the interface for the \textit{Digital / Technological Transformation} context. The upper panel lets analysts select an atomic rule and inspect its membership function together with the cohort distribution of the underlying feature; a \textit{Numerical Audit} table on the right reports, for each case, the crisp input value and the corresponding membership degree $\mu$. The lower panel, \textit{Derived Rules Overview}, summarizes satisfaction in the composite rules and the resulting context-alignment score $C_i$ for the full cohort.  

\begin{figure}[H]
\centering
\includegraphics[width=\textwidth]{figures/context.png}
\caption{\textit{Context} module for the \textit{Digital / Technological Transformation} run, combining an atomic-rule selector, the membership-function chart with case distribution, and the numerical audit table for \textit{Age investment window} (trapezoidal), together with a derived-rules overview reporting membership degrees and $C_i$ for cases ordered by decreasing context alignment.}
\label{fig:context}
\end{figure}

In the exported view, the analyst has selected the trapezoidal atomic rule \textit{Age investment window}, which defines the acceptable age range for developmental investment during the transformation programme (Appendix~\ref{app:context_specs}). The chart links each employee's age to its membership $\mu$, while the audit table makes the same mapping explicit case by case. Below, \textit{Derived Rules Overview} exposes the degrees of satisfaction for \textit{Upskilling capacity}, \textit{Career scalability}, and \textit{Investment horizon}, together with the aggregated score $C_i$. The table is sorted by $C_i$ in descending order, so \textit{Noah Lewis} ($C_i=0.7100$) appears first and \textit{Evelyn Taylor} ($C_i=0.7000$) second, even though Evelyn leads the prioritization ranking in Fig.~\ref{fig:overview} because of higher cooperative risk. Across the cohort, $C_i$ ranged from 0.1500 to 0.7100 in this run. This view therefore complements the global prioritization summary by making the contextual evidence auditable rule by rule and case by case.

\subsubsection{\textit{Explain}: local case interpretation}

For a selected case, the \textit{Explain} module assembles the four explanation levels defined in Section~\ref{sec:architecture_improvements} into a single audit view. Figure~\ref{fig:explain} exports the interface for \textit{Evelyn Taylor}, the highest-priority case in Fig.~\ref{fig:overview}. Fig.~\ref{fig:explain}-A) shows the case selector, the deterministic natural-language summary, and the headline statistics on which that summary rests: $P_i$, $R_i$, $C_i$, and $\lambda$, together with the additive decomposition in Equation~\eqref{eq:priority_decomposition}. Fig.~\ref{fig:explain}-B) reports the aggregated ML attributions for the ten most influential features and, below, the fuzzy memberships that trace how strongly each contextual rule is satisfied for the selected case.

\begin{figure}[t]
\begin{adjustwidth}{-0.8\extralength}{0cm}
	\centering
\includegraphics[width=1.3\textwidth]{figures/explain_high_prioritization.pdf}
\end{adjustwidth}
\caption{\textit{Explain} module for \textit{Evelyn Taylor} at $\lambda=0.5$ in the \textit{Digital / Technological Transformation} context, with case selection, natural-language explanation, summary statistics, aggregated bar charts of the ten most influential ML features, and a fuzzy context traceability panel reporting membership degrees for the policy rules.}
\label{fig:explain}
\end{figure}

Two cases from the \textit{Digital / Technological Transformation} run illustrate the interpretation produced by the \textit{Explain} module at $\lambda=0.5$. The first is \textit{Evelyn Taylor}, ranked first in Fig.~\ref{fig:overview} and reproduced in Figure~\ref{fig:explain}; the second is \textit{Noah Lewis}, who ranks low in the same overview despite having the highest $C_i$ in the cohort. For space reasons, only the \textit{Evelyn Taylor} view is shown in the figure; the contrasting \textit{Noah Lewis} summary is reported below in textual form. In both cases, the summaries were reproduced from CADEMAS-ML and follow its deterministic template: the final prioritization score $P_i$ is reported first, then $\lambda$, the levels of $R_i$ and $C_i$, the main ML drivers, the contextual rules that support or constrain priority, and a tier-based recommendation derived from $P_i$.

\textit{Evelyn Taylor} combines critical predictive risk ($R_i=0.9555$) with high contextual alignment ($C_i=0.7000$), yielding $P_i(0.5)=0.8277$ and a \textit{Tier 1} recommendation. Fig.~\ref{fig:explain}-A) displays these values together with the natural-language synthesis; Fig.~\ref{fig:explain}-B) highlights Department (Sales) and Overtime (Yes) among the strongest ML drivers and shows high membership on Relevant years at company and Age investment window in the fuzzy traceability panel:

\begin{quote}
\small
The case prioritization score for \textit{Evelyn Taylor} is 83\%. This decision reflects a balanced posture ($\lambda=0.50$), combining critical predictive risk (96\%) with high contextual alignment (70\%). On the predictive side, risk was driven mainly by Department (Sales) and Overtime (Yes). In context, the case aligns strongly with the policy rules, particularly on Relevant years at company and Age investment window.

Suggestion [\textit{Tier 1} -- High Priority]: Recommended to proceed with intervention or resource allocation immediately.
\end{quote}

\textit{Noah Lewis} provides a contrasting interpretation. Although this case has the highest contextual alignment in the cohort ($C_i=0.7100$), its cooperative risk is low ($R_i=0.0651$). The resulting prioritization score is therefore $P_i(0.5)=0.3875$, and the module assigns \textit{Tier 3}:

\begin{quote}
\small
The case prioritization score for \textit{Noah Lewis} is 39\%. This decision reflects a balanced posture ($\lambda=0.50$), combining low predictive risk (7\%) with high contextual alignment (71\%). On the predictive side, risk was driven mainly by Years with current manager (8) and Number of companies worked (1). In context, the case aligns strongly with the policy rules, particularly on Relevant years at company and Age investment window.

Suggestion [\textit{Tier 3} -- Low Priority]: Intervention or resource allocation is not justified under current prediction and context guidelines.
\end{quote}

In practical terms, CADEMAS-ML explains not only why a case is prioritized, but also why another case with strong contextual alignment may still receive a low final score when predictive risk remains limited. A third illustrative pattern is provided by \textit{Theodore Adams} ($R_i=0.9803$, $C_i=0.4400$, $P_i(0.5)=0.7101$), who receives \textit{Tier 2} because high predictive risk and only moderate contextual alignment produce borderline evidence under the current policy.

\subsubsection{\textit{Robustness}: sensitivity of the priority ranking}

The \textit{Robustness} module implements the sensitivity analysis defined in Section~\ref{sec:architecture_improvements}. Figure~\ref{fig:robustness} illustrates the exported Altair charts for the \textit{Digital / Technological Transformation} run using $\Lambda=\{0,0.25,0.5,0.75,1\}$ and the ten highest-priority cases at $\lambda=0.5$ (the head of the ranking in Fig.~\ref{fig:overview}).

Fig.~\ref{fig:robustness}-A) shows the bump chart of priority ranks across $\lambda$. Fig.~\ref{fig:robustness}-B) reports the ranks box plot, which summarizes the rank distribution of each displayed case over the same sweep; cases are ordered from left to right by median rank and, in case of ties, by interquartile range. Fig.~\ref{fig:robustness}-C) shows the rank acceptability indices, estimating the relative frequency with which each case occupies each rank position over $\Lambda$. 

When $\lambda$ approaches 0 and contextual alignment receives greater weight, cases outside the top ten at $\lambda=0.5$, such as \textit{Noah Lewis} and \textit{Emma Smith}, rise in Fig.~\ref{fig:robustness}-A) because they combine high $C_i$ with low $R_i$. When $\lambda$ approaches 1, the high-risk cases at the head of Fig.~\ref{fig:overview}---such as \textit{Lucas Wright} and \textit{Theodore Adams}---dominate the ranking regardless of transformation alignment. \textit{Evelyn Taylor} remains near the top across intermediate values of $\lambda$ because both $R_i$ and $C_i$ are high. The box plot in Fig.~\ref{fig:robustness}-B) makes this contrast visible at a glance: stable high-priority cases exhibit compact rank distributions, whereas policy-sensitive cases show wider spreads. Fig.~\ref{fig:robustness}-C) complements both views by quantifying how often each case attains each rank over the $\lambda$ grid.

\begin{figure}[H]
\centering
\includegraphics[width=\textwidth]{figures/robustness.pdf}
\caption{Robustness analysis for the \textit{Digital / Technological Transformation} context exported from CADEMAS-ML over $\Lambda=\{0,0.25,0.5,0.75,1\}$ for the ten highest-priority cases at $\lambda=0.5$, comprising a bump chart of priority ranks across $\lambda$, a ranks box plot of rank distributions over the same sweep, and rank acceptability indices.}
\label{fig:robustness}
\end{figure}

\section{Conclusion and future works}\label{sec:conclusions}

This paper presented CADEMAS-ML, a software implementation that operationalizes the CADEMAS framework for cooperative, context-aware, ML-based prioritization. The tool preserves the conceptual separation proposed by CADEMAS between ADM-unit outputs, integration, contextual evaluation, and final decision construction, while providing an executable workflow based on H2O MOJO models, configurable model weights, fuzzy context JSON files, and an adjustable policy parameter $\lambda$.

The main contribution of CADEMAS-ML is not the introduction of a new attrition model, but the transformation of the CADEMAS architecture into an auditable decision-support tool. In this respect, two modules are especially relevant. The \textit{Robustness} module extends the static prioritization workflow by evaluating how rankings change under alternative balances between predictive risk and contextual alignment. The \textit{Explain} module adds local interpretability by decomposing the priority score, estimating model-agnostic feature attributions, tracing fuzzy-rule memberships, and producing natural-language recommendations. These capabilities were not present in the original CADEMAS framework and were not implemented in the attrition instantiation of \citet{Godz2026}.

The employee-attrition case study illustrates how the same set of predictive models can lead to different prioritization outcomes depending on the contextual policy adopted. It also shows how individual cases can be inspected before action is taken. Future work will focus on broadening the supported ADM-unit types beyond predictive MOJO models, incorporating additional explanation methods when model internals are available, improving governance support for context-versioning and audit trails, and validating the tool in other prioritization domains such as fraud triage, credit review, resource allocation, and medical or social-service prioritization.

\reftitle{References}
\bibliography{biblio.bib}

\newpage
\appendixtitles{yes}
\appendixstart
\appendix
\section{Attrition example dataset and model features}\label{app:attrition_features}

The 20-case attrition example distributed with CADEMAS-ML contains 36 columns: \texttt{Case\_ID}, 34 predictor variables, and the \texttt{attrition} outcome label. The predictors cover demographic, professional, compensation, performance, and organizational attributes from the \textit{IBM HR Analytics Employee Attrition and Performance} Kaggle dataset used in \citet{Godz2026}. Five variables (\texttt{age}, \texttt{gender}, \texttt{department}, \texttt{job\_role}, and \texttt{job\_level}) are shared by all four MOJO models. Two additional columns (\texttt{employee\_count} and \texttt{employee\_number}) are retained in the CSV file but are not used by any model in the example bundle. Each MOJO receives the full case matrix and selects its required variables according to the schema stored in the model-definition JSON file.

Table~\ref{tab:attrition_feature_assignment} summarizes the feature assignment. The Human Resources model reads \texttt{over18} from the CSV file; the corresponding field is named \texttt{over\_18} in the model metadata.

\begin{table}[H]
\caption{Predictor variables in the attrition example and the MOJO models that use each one. $M_1$--$M_4$ denote Culture and Wellbeing, Finance, Operations and Development, and Human Resources, respectively.\label{tab:attrition_feature_assignment}}
\small
\begin{adjustwidth}{-0.3\extralength}{0cm}
\begin{tabularx}{1.1\textwidth}{l|cccc}
\toprule
\textbf{Predictor} & \shortstack[c]{$M_1$\\{\scriptsize (Culture and Wellbeing)}} & \shortstack[c]{$M_2$\\{\scriptsize (Finance)}} & \shortstack[c]{$M_3$\\{\scriptsize (Operations and Development)}} & \shortstack[c]{$M_4$\\{\scriptsize (Human Resources)}} \\
\midrule
age & $\checkmark$ & $\checkmark$ & $\checkmark$ & $\checkmark$ \\
business\_travel & $\checkmark$ &  &  &  \\
daily\_rate &  & $\checkmark$ &  &  \\
department & $\checkmark$ & $\checkmark$ & $\checkmark$ & $\checkmark$ \\
distance\_from\_home &  &  &  & $\checkmark$ \\
education &  &  &  & $\checkmark$ \\
education\_field &  &  &  & $\checkmark$ \\
employee\_count &  &  &  &  \\
employee\_number &  &  &  &  \\
environment\_satisfaction & $\checkmark$ &  &  &  \\
gender & $\checkmark$ & $\checkmark$ & $\checkmark$ & $\checkmark$ \\
hourly\_rate &  & $\checkmark$ &  &  \\
job\_involvement & $\checkmark$ &  &  &  \\
job\_level & $\checkmark$ & $\checkmark$ & $\checkmark$ & $\checkmark$ \\
job\_role & $\checkmark$ & $\checkmark$ & $\checkmark$ & $\checkmark$ \\
job\_satisfaction & $\checkmark$ &  &  &  \\
marital\_status &  &  &  & $\checkmark$ \\
monthly\_income &  & $\checkmark$ &  &  \\
monthly\_rate &  & $\checkmark$ &  &  \\
num\_companies\_worked &  &  & $\checkmark$ &  \\
over18 &  &  &  & $\checkmark$ \\
over\_time & $\checkmark$ &  &  &  \\
percent\_salary\_hike &  & $\checkmark$ &  &  \\
performance\_rating &  &  & $\checkmark$ &  \\
relationship\_satisfaction & $\checkmark$ &  &  &  \\
standard\_hours &  &  &  & $\checkmark$ \\
stock\_option\_level &  & $\checkmark$ &  &  \\
total\_working\_years &  &  & $\checkmark$ &  \\
training\_times\_last\_year &  &  & $\checkmark$ &  \\
work\_life\_balance & $\checkmark$ &  &  &  \\
years\_at\_company &  &  & $\checkmark$ &  \\
years\_in\_current\_role &  &  & $\checkmark$ &  \\
years\_since\_last\_promotion &  &  & $\checkmark$ &  \\
years\_with\_curr\_manager &  &  & $\checkmark$ &  \\
\midrule
\textbf{Total features} & 12 & 11 & 13 & 11 \\
\bottomrule
\end{tabularx}
\end{adjustwidth}
\end{table}

\section{Context specification for the \textit{Digital / Technological Transformation} scenario}\label{app:context_specs}

The example repository includes two context configurations, but the case study focuses on \textit{Digital / Technological Transformation}. This appendix reports the corresponding fuzzy specification in closed form so that every membership degree and the final context-alignment score $C_i$ can be audited without inspecting application code. The rules instantiate the generic formalism of Equations~\eqref{eq:atomic_mu}--\eqref{eq:context_score}; equivalent parameters are stored in the repository JSON file consumed by CADEMAS-ML\footnote{\url{https://github.com/pnovoa/cademas-app/blob/main/example_attrition/context/context_digital_transformation.json}.}. They are illustrative policy encodings, not hard-coded assumptions of the tool.

\subsection{Generic membership functions}\label{app:context_math_mf}

Each numerical atomic rule is evaluated through one of four parametric membership functions. For an input value $x$, the \emph{linear increasing} function with breakpoints $(a,b)$, $a<b$, is
\begin{equation}
\mu^{\uparrow}_{a,b}(x)=
\begin{cases}
0, & x \le a,\\[2pt]
\dfrac{x-a}{b-a}, & a < x < b,\\[6pt]
1, & x \ge b,
\end{cases}
\label{eq:mf_linc}
\end{equation}
and the \emph{linear decreasing} function with breakpoints $(a,b)$ is its mirror image
\begin{equation}
\mu^{\downarrow}_{a,b}(x)=1-\mu^{\uparrow}_{a,b}(x)=
\begin{cases}
1, & x \le a,\\[2pt]
\dfrac{b-x}{b-a}, & a < x < b,\\[6pt]
0, & x \ge b.
\end{cases}
\label{eq:mf_ldec}
\end{equation}
The \emph{triangular} function with support $(a,c)$ and peak $b$, $a<b<c$, is
\begin{equation}
\Lambda_{a,b,c}(x)=
\begin{cases}
\dfrac{x-a}{b-a}, & a < x \le b,\\[6pt]
\dfrac{c-x}{c-b}, & b < x < c,\\[6pt]
0, & \text{otherwise},
\end{cases}
\label{eq:mf_tri}
\end{equation}
and the \emph{trapezoidal} function with feet $(a,d)$ and shoulders $(b,c)$, $a<b\le c<d$, is
\begin{equation}
\Pi_{a,b,c,d}(x)=
\begin{cases}
\dfrac{x-a}{b-a}, & a < x < b,\\[6pt]
1, & b \le x \le c,\\[2pt]
\dfrac{d-x}{d-c}, & c < x < d,\\[6pt]
0, & \text{otherwise}.
\end{cases}
\label{eq:mf_trap}
\end{equation}
Categorical attributes are handled by a direct map $\kappa:\mathcal{S}\to[0,1]$ that assigns a degree to each symbolic value $s\in\mathcal{S}$, with unspecified categories defaulting to $0$.

\subsection{Atomic rules}\label{app:context_math_atomic}

Seven atomic rules evaluate recent training intensity, education level, company tenure, tenure in the current role, age in the investment window, accumulated experience, and role adaptability to technological change. For case $x_i$, the corresponding memberships are
\begin{align}
\mu_{\textsf{train}}(x_i) &= \mu^{\uparrow}_{1,6}\!\left(\textsf{training\_times\_last\_year}_i\right), \label{eq:atom_train}\\
\mu_{\textsf{edu}}(x_i) &= \mu^{\uparrow}_{2,5}\!\left(\textsf{education}_i\right), \label{eq:atom_edu}\\
\mu_{\textsf{yrsCo}}(x_i) &= \Lambda_{3,10,20}\!\left(\textsf{years\_at\_company}_i\right), \label{eq:atom_yrsco}\\
\mu_{\textsf{yrsRole}}(x_i) &= \Lambda_{1,6,11}\!\left(\textsf{years\_in\_current\_role}_i\right), \label{eq:atom_yrsrole}\\
\mu_{\textsf{age}}(x_i) &= \Pi_{22,28,50,57}\!\left(\textsf{age}_i\right), \label{eq:atom_age}\\
\mu_{\textsf{exp}}(x_i) &= \mu^{\uparrow}_{2,10}\!\left(\textsf{total\_working\_years}_i\right), \label{eq:atom_exp}\\
\mu_{\textsf{role}}(x_i) &= \kappa_{\textsf{role}}\!\left(\textsf{job\_role}_i\right), \label{eq:atom_role}
\end{align}
where the role-adaptability map assigns expert scores to each job role:
\begin{equation}
\kappa_{\textsf{role}}(s)=
\begin{cases}
0.9, & s\in\{\text{Research Scientist},\ \text{Research Director}\},\\
0.8, & s=\text{Manager},\\
0.7, & s\in\{\text{Manufacturing Director},\ \text{Healthcare Representative}\},\\
0.6, & s=\text{Sales Executive},\\
0.5, & s=\text{Human Resources},\\
0.4, & s\in\{\text{Laboratory Technician},\ \text{Sales Representative}\}.
\end{cases}
\label{eq:kappa_role}
\end{equation}
Each expression can be read directly: for instance, Equation~\eqref{eq:atom_age} states that ages between 28 and 50 fully satisfy the investment window ($\mu=1$), the degree decays linearly toward the feet at 22 and 57, and ages outside $[22,57]$ are excluded.

\subsection{Derived rules and final aggregation}\label{app:context_math_derived}

The aggregation operators of Equation~\eqref{eq:derived_mu}, applied to a vector $\mathbf{u}=(u_1,\ldots,u_q)$, are
\begin{equation}
\textsc{and}(\mathbf{u})=\min_{j} u_j,\quad
\textsc{or}(\mathbf{u})=\max_{j} u_j,\quad
\textsc{prod}(\mathbf{u})=\prod_{j=1}^{q} u_j,\quad
\textsc{avg}(\mathbf{u})=\frac{1}{q}\sum_{j=1}^{q} u_j.
\label{eq:agg_ops}
\end{equation}
With these operators, three derived criteria combine the atomic memberships: \emph{upskilling capacity} (training and education), \emph{career scalability} (company and role tenure), and \emph{investment horizon} (age window and experience). Formally,
\begin{align}
\mu_{\textsf{upskill}}(x_i) &= \textsc{avg}\!\left(\mu_{\textsf{train}}(x_i),\,\mu_{\textsf{edu}}(x_i)\right)
= \tfrac{1}{2}\!\left(\mu_{\textsf{train}}(x_i)+\mu_{\textsf{edu}}(x_i)\right), \label{eq:der_upskill}\\
\mu_{\textsf{career}}(x_i) &= \textsc{prod}\!\left(\mu_{\textsf{yrsCo}}(x_i),\,\mu_{\textsf{yrsRole}}(x_i)\right)
= \mu_{\textsf{yrsCo}}(x_i)\,\mu_{\textsf{yrsRole}}(x_i), \label{eq:der_career}\\
\mu_{\textsf{invest}}(x_i) &= \textsc{prod}\!\left(\mu_{\textsf{age}}(x_i),\,\mu_{\textsf{exp}}(x_i)\right)
= \mu_{\textsf{age}}(x_i)\,\mu_{\textsf{exp}}(x_i). \label{eq:der_invest}
\end{align}
The choice of operator carries a clear semantics: \textsc{avg} lets strong training compensate for weaker formal education in \emph{upskilling capacity}, whereas the \textsc{prod} operator in \emph{career scalability} and \emph{investment horizon} is conjunctive, so a near-zero membership in either input collapses the derived degree toward zero. Finally, the context-alignment score is the average-based soft conjunction $H$ over the upskilling, role-adaptability, career-scalability, and investment-horizon memberships:
\begin{equation}
C_i = \textsc{avg}\!\left(\mu_{\textsf{upskill}}(x_i),\,\mu_{\textsf{role}}(x_i),\,\mu_{\textsf{career}}(x_i),\,\mu_{\textsf{invest}}(x_i)\right)
= \frac{1}{4}\sum_{r\in\mathcal{R}} \mu_{r}(x_i),
\label{eq:context_score_digital}
\end{equation}
with $\mathcal{R}=\{\textsf{upskill},\textsf{role},\textsf{career},\textsf{invest}\}$. Equation~\eqref{eq:context_score_digital} is the explicit instance of the generic policy in Equation~\eqref{eq:context_score} for this context and defines the contextual evidence surfaced in the \textit{Context} module (Figure~\ref{fig:context}) and traced at case level in the \textit{Explain} module.


%\PublishersNote{}

\end{document}
